% NOTA: este fichero se incluye desde sp1.tex tras \appendix.
% ==========================================================================
% SP1 · material desplazado fuera del cuerpo del capítulo
%
% Se compila como apendice de sp1.tex mediante \input. Recoge, sin perdida,
% el material que la pasada editorial final retiro del cuerpo. Al integrar el
% capitulo en el TFM cada bloque puede reubicarse en su destino definitivo:
%
%   A. Protocolo Monte Carlo pareado      -> capítulo de metodología
%   B. Pseudocódigo de CBBA-RB y GRAPE    -> anexo de algoritmos
%   C. Taxonomía F-I--F-IV y su tabla     -> anexo de métodos
%   D. Glosario de la familia F-I          -> notación
%   E. Trazabilidad e hipótesis H1/H5     -> conclusiones
%   F. Nota de congelación de semillas    -> anexo de reproducibilidad
%
% Cada bloque conserva el texto exacto que tenía en el cuerpo.
% ==========================================================================

% --------------------------------------------------------------------------
% A. Protocolo Monte Carlo pareado (procedía de SP1, p. 4)

\anxsection{Pseudocódigo de los métodos de N3}{anx:pseudocodigo}
% --------------------------------------------------------------------------
\begin{minipage}[t]{.49\linewidth}
{\scriptsize\textbf{Algoritmo 1. Capacity-CBBA-RB.}\par
\begin{algorithmic}[1]
\Require tabla local $R_i$, barreras $h_i[k]$, versión $v_i$
\State recibir $(k,b,C,v,\tau)$ y descartar versiones antiguas
\State fusionar ganadores por $(\Delta D,-d,-\tau)$
\State reconstruir el prefijo de cada $k$ hasta cubrir $m_k$
\If{$i$ queda fuera del prefijo de $k$}
  \State $h_i[k]\gets\max\{h_i[k],b_{ik}\}$; liberar $k$
\EndIf
\State calcular $r_{ik}=[m_k-\sum_{j\in\widehat C_{ik}}c_j]_+$
\If{$i$ está libre}
  \State elegir $k^\star=\arg\max_k(\min\{c_i,r_{ik}\},-d_{ik},-\tau_i)$
  \State pujar solo si $b_{ik^\star}>h_i[k^\star]$; elevar $v_i$
\Else
  \State conservar el registro vigente
\EndIf
\State difundir cambios; refresco completo cada $N$ activaciones
\Ensure coalición candidata, versiones y bytes emitidos
\end{algorithmic}}
\end{minipage}\hfill
\begin{minipage}[t]{.49\linewidth}
{\scriptsize\textbf{Algoritmo 2. Weighted-/Pair-GRAPE.}\par
\begin{algorithmic}[1]
\Require perfil común $a$, fase $q\in\{0,1\}$, contador local
\State intercambiar $(c_i,d_i[\cdot],v_i)$ y verificar el hash de $a$
\State calcular la mejor desviación que reduzca $(D,J)$
\State proponer $(\Delta D,\Delta J,i,a_i',v_i)$ o propuesta nula
\For{$\ell=1,\ldots,N-2$}
  \State inundar el máximo por orden total; reintentar hasta 4 veces
\EndFor
\State en $\ell=N-1$, aplicar una única propuesta en todos los nodos
\State incrementar época y difundir versión y hash del perfil
\If{no hay propuesta y $q=0$ en Pair-GRAPE}
  \State fijar $q\gets1$; enumerar pares con un vecino
\ElsIf{no hay propuesta}
  \State terminar y verificar coincidencia de perfiles
\Else
  \State volver a la etapa de propuesta
\EndIf
\Ensure perfil terminal común o fallo operacional por desacuerdo
\end{algorithmic}}
\end{minipage}

% --------------------------------------------------------------------------
% C. Taxonomía de familias decisionales de N4 (procedía de SP1, p. 13)

\anxsection{Atlas de reglas y de orden de revisión}{anx:reglas}
% ------------------------------------------------------------------
\begin{figure}[H]\centering
\tikzpanelfit{.98\linewidth}{8.4cm}{figures/n4/n4_order_revision.tex}
{Orden estratégico y regla de revisión}
\caption{Vecindades $\mathcal V_1$, $\mathcal V_2$ y $\mathcal V_3$ frente a las
reglas BR, ASR y LLL.}
\viuderivedsource{la revisión de \citet{sandholm2010population} y el
aprendizaje log-lineal de \citet{mardenShamma2012LogLinear}}
\end{figure}

\begin{figure}[H]\centering
\tikzpanelfit{.98\linewidth}{11.2cm}{figures/n4/n4_algorithm_atlas_fi.tex}
{BR, Geo-ASR y Geo-LLL}
\caption{BR, Geo-ASR y Geo-LLL sobre un mismo estado inicial: regla de selección y
desviación aceptada dentro de $\mathcal V_1(a)$.}
\viuderivedsource{la revisión de \citet{sandholm2010population} y el
aprendizaje log-lineal de \citet{mardenShamma2012LogLinear}}
\end{figure}

\begin{figure}[H]\centering
\tikzpanelfit{.98\linewidth}{11.2cm}{figures/n4/n4_algorithm_atlas_coalitional.tex}
{2BR, C3 y DMIS+TX}
\caption{2BR y C3 sobre el mismo estado: desviaciones de orden dos y tres. DMIS+TX:
grafo de conflictos, conjunto independiente y confirmación transaccional con
versiones.}
\viuderivedsource{\citet{luby1986MIS} y \citet{grayLamport2006Commit}}
\end{figure}

% ------------------------------------------------------------------
% J. Esquema de los comparadores continuos (procedia de SP1, F-II/F-III)

\anxsection{Protocolo experimental y análisis de escala}{anx:escala}
%    Destino: capitulo comun de metodologia.
% ------------------------------------------------------------------
\subsection{Mundos, generadores y contraste pareado}

Los mundos se generan con cinco geometrías normalizadas
(Figura~\ref{fig:geometrias}) y la Tabla~\ref{tab:protocolo} recoge sus
parámetros junto con la regla de inferencia. Cada semilla genera un mundo
independiente. Dentro de un mismo mundo todos los
métodos reciben los mismos AMR, cargas y parámetros, y el \emph{bootstrap}
remuestrea mundos completos por celda. Las ejecuciones fallidas o censuradas se
conservan en los datos, y Holm corrige cada familia de contrastes. Las campañas
siguen el protocolo Monte Carlo pareado definido en el capítulo de metodología.

\begin{figure}[H]\centering
\tikzpanel{.98\linewidth}{9.8cm}{figures/compact/fig03_scenarios.tex}
{Cinco geometrías y retirada de un AMR}
\caption{Cinco generadores espaciales normalizados y el tratamiento con retirada
de un AMR.}
\label{fig:geometrias}
\viuownsource
\end{figure}

\begin{table}[H]
\begin{minipage}[t]{.57\linewidth}
{\footnotesize
\textbf{(a) Generadores espaciales.}\par
\renewcommand{\arraystretch}{1.08}
\begin{tabularx}{\linewidth}{@{}p{1.55cm}p{3.25cm}X@{}}
\toprule
\textbf{Geo.}&\textbf{AMR}&\textbf{Cargas}\\
\midrule
Aleatorio&$U([0,L_x]\!\times\![0,L_y])$&igual\\
Agrupado&uniforme&dos centros + $\mathcal N(0,\Sigma)$\\
Separado&$x\le.35L_x$&$x\ge.65L_x$\\
Anillo&$r\in[.36,.48]L$&$r\in[.04,.22]L$\\
Pasillo&$\sigma_y=.04L_y$&$\sigma_y=.08L_y$\\
\bottomrule
\end{tabularx}
\par\smallskip
Agrupado, Anillo y Pasillo aplican \texttt{clip}; la retirada de un AMR conserva
la geometría Aleatorio.
}
\end{minipage}\hfill
\begin{minipage}[t]{.39\linewidth}
{\footnotesize
\textbf{(b) Regla estadística común.}\par
\renewcommand{\arraystretch}{1.08}
\begin{tabularx}{\linewidth}{@{}>{\raggedright\arraybackslash}p{1.9cm}>{\raggedright\arraybackslash}X@{}}
\toprule
\textbf{Respuesta}&\textbf{Inferencia}\\
\midrule
Binaria&riesgo pareado; \mbox{McNemar}\\
Continua&mediana + IC \emph{bootstrap}; signo\\
Tres o más&Friedman + Kendall $W$\\
Censura&tiempo truncado; \mbox{supervivencia}\\
\bottomrule
\end{tabularx}%
}%
\end{minipage}
\caption{(a) Parámetros de los cinco generadores espaciales. (b) Regla de
inferencia aplicada según el tipo de respuesta.}
\label{tab:protocolo}
\viuownsource
\end{table}

% ------------------------------------------------------------------
% M. Sensibilidad de P(h_c* > 3) (procedia del cuerpo, N4.E9)
% ------------------------------------------------------------------

\anxsection{Sensibilidad del primer escape conectado}{anx:sensibilidad}

Esos \NFourHStarGtThreeN{} casos se concentran en presión alta y capacidad
homogénea. Dentro del bloque confirmatorio, que fija $N/K=\NFourMainBlockNk$, la
proporción sin escape detectado fue del \NFourNoEscRhoHigh\,\% con
$\rho=0{,}85$ y del \NFourNoEscRhoLow\,\% con $\rho=0{,}70$
($n=\NFourNoEscRhoHighN$ por nivel), y del \NFourNoEscCvZero\,\% con
$\mathrm{CV}=0$, mientras que en los tres niveles de dispersión
restantes no pasó del \NFourNoEscCvMid\,\%
($n=\NFourNoEscCvZeroN$ por nivel).

La relación $N/K$ se barre en un bloque de sensibilidad aparte, de
\NFourRatioBlockN{} mundos, así que sus cifras no comparten denominador con las
anteriores: \NFourNoEscNkTwo\,\% con $N/K=2$, \NFourNoEscNkThree\,\% con $N/K=3$ y
\NFourNoEscNkFour\,\% con $N/K=4$ ($n=\NFourNoEscNkTwoN$ por nivel). Dentro de
ese bloque, la tasa sin escape disminuye de forma marcada al pasar de dos a cuatro
AMR por carga; no se compara con CV ni con la presión, que provienen de otro
diseño.

Los dos bloques responden a diseños distintos: el confirmatorio fija
$N/K=\NFourMainBlockNk$ y barre CV y presión, mientras que el de sensibilidad
barre $N/K$. Sus porcentajes no son términos de una misma descomposición.

% ------------------------------------------------------------------
% N. Formulacion de los comparadores continuos (procedia del cuerpo)
% ------------------------------------------------------------------

\anxsection{Formulación de los comparadores continuos}{anx:formulacion}

\begin{minipage}[t]{.49\linewidth}
\textbf{Dinámicas poblacionales}
\citep{sandholm2010population,quijano2017population}.
Cada AMR reparte una unidad de masa entre las $K+1$ opciones, con $x_i$ en el
símplex $\mathcal X_i$. La cobertura es $Q_k(x)=\sum_ic_i^{\rm pay}x_{ik}$, el
déficit relativo $r_k=[1-Q_k(x)/m_k]_+$ y $\overline{d}_{ik}$ la distancia
normalizada. Con ellos:
\begin{equation}
\begin{aligned}
\label{eq:n4-fii-potential}
\widetilde\Phi(x)&=-\frac{\alpha}{2}\sum_kr_k(x)^2
-\frac{\beta}{N}\sum_{i,k}\overline{d}_{ik}x_{ik},\\[-1mm]
F_{ik}&=\partial_{x_{ik}}\widetilde\Phi .
\end{aligned}
\tag{7}
\end{equation}
\begin{equation}
\begin{aligned}
\dot x_{ik}^{\rm Rep}&=x_{ik}(F_{ik}-\overline{F}_i),\\
\dot x_{ik}^{\rm Smith}&=\sum_qx_{iq}[F_{ik}-F_{iq}]_+
\\[-1mm]&\quad-x_{ik}\sum_q[F_{iq}-F_{ik}]_+,\\
\dot x_{ik}^{\rm BNN}&=[F_{ik}-\overline{F}_i]_+
\\[-1mm]&\quad-x_{ik}\sum_q[F_{iq}-\overline{F}_i]_+,\\
\dot x_{ik}^{\rm Logit}&=
\frac{e^{F_{ik}/\tau}}{\sum_qe^{F_{iq}/\tau}}-x_{ik}.
\end{aligned}
\tag{8}
\end{equation}
$\overline{F}_i=\sum_qx_{iq}F_{iq}$ es el pago medio de $i$ y $\overline{d}_{ik}=d_{ik}/\max_{j,q}d_{jq}$.
\end{minipage}\hfill
\begin{minipage}[t]{.49\linewidth}
\textbf{Primal-dual con restricción compartida.}
\begin{equation}
\begin{aligned}
g_k(x)&=1-\frac{Q_k(x)}{m_k}\le0,\\[-1mm]
J_i&=\frac1N\sum_k\overline{d}_{ik}x_{ik}+\frac\epsilon2\|x_i\|^2 .
\end{aligned}
\tag{9}
\end{equation}
\begin{equation}
\begin{aligned}
\min_{x_i\in\mathcal X_i}\quad&J_i(x_i)\\[-1mm]
\mathrm{s.a.}\quad&g(x)\le0,\qquad i=1,\ldots,N .
\end{aligned}
\tag{10}
\end{equation}
\begin{equation}
\begin{aligned}
\dot x_i&=\Pi_{T_{\mathcal X_i}}[-\nabla_{x_i}\mathcal L_i],\\[-1mm]
\dot\lambda_k&=\Pi_{T_{\mathbb R_+}}[g_k(x)] .
\end{aligned}
\tag{11}
\label{eq:n4-primal-dual}
\end{equation}
con $\mathcal L_i=J_i(x_i)+\lambda^{\!\top}g(x)$.
\begin{equation}
\begin{aligned}
\dot\lambda_{ik}&=\Pi_{T_{\mathbb R_+}}
\!\bigl[\widehat g_{ik}\\[-1mm]
&\qquad-\gamma\!\sum_{j\in\mathcal N_i}
(\lambda_{ik}-\lambda_{jk})\bigr].
\end{aligned}
\tag{12}
\label{eq:n4-consensus}
\end{equation}
$\widehat g_{ik}$ es la estimación local que $i$ mantiene de $g_k$.
\begin{equation}
\label{eq:n4-kkt}
\begin{gathered}
0\in\mathbf F(x^\star)+\nabla g(x^\star)^\top\lambda^\star
+N_{\mathcal X}(x^\star),\\
g(x^\star)\le0,\ \lambda^\star\ge0,\ 
{\lambda^\star}^{\!\top}g(x^\star)=0 .
\end{gathered}
\tag{13}
\end{equation}
$\mathbf F(x)=(\nabla_{x_i}J_i(x))_{i}$ es el pseudogradiente y $N_{\mathcal X}$ el
cono normal a $\mathcal X=\prod_i\mathcal X_i$.
\end{minipage}

% ------------------------------------------------------------------
% O. Nota de interfaz SP1 -> SP2 (procedia de la sintesis)
% ------------------------------------------------------------------

\anxsection{Orden conectado frente a mejora}{anx:hstar}
% ------------------------------------------------------------------
\begin{figure}[H]\centering
\pdfpanel{.90\linewidth}{5.2cm}{assets/figures/n4-v4/n4_hstar_gap_link.pdf}
{Orden conectado frente a mejora}
\caption{Mejora de brecha por categoría de $h_c^\star$, sobre los
\NFourHStarLinkN{} mundos en que BR terminó con salida factible y las tres brechas
quedan definidas. Con $h_c^\star=3$ ($n=\NFourHStarLinkThreeN$) y sin escape
detectado ($n=\NFourHStarLinkGtThreeN$) se muestran las observaciones
individuales, no una densidad estimada.}
\viuownsource
\end{figure}

% ------------------------------------------------------------------
% L. Protocolo experimental de SP1 (procedia del cuerpo, p. 3)

% ------------------------------------------------------------------
% destino: documentación de interfaz SP1--SP2
% ------------------------------------------------------------------


\anxsection{Diagnóstico post hoc de conectividad}{anx:conectividad}

Como diagnóstico posterior a la campaña se comprobó si una medida agregada de
densidad basta para ordenar los regímenes. Sobre el área de trabajo $A=L_xL_y$,
la escala de conectividad de los grafos geométricos aleatorios
\citep{penrose2003RandomGeometric} sugiere el índice adimensional
$\Gamma_R=\pi NR^2/(A\log N)$, que aquí no es un umbral, y ese índice no separa
los regímenes. Mundos etiquetados igual ocupan
rangos solapados, porque cuatro de los cinco generadores no reparten los AMR de
forma uniforme y la geometría cambia la densidad efectiva a igualdad de
multiplicador.

La conectividad algebraica tampoco los separa mundo a mundo. Con $L_G$ el
laplaciano del grafo de comunicación, $\lambda_2(L_G)$ ---su segundo autovalor más
pequeño--- crece cuanto más fuertemente unida está la red. Decrece de forma
monótona en mediana al ralear el grafo, pero los rangos de dos regímenes
contiguos se solapan igual que los de $\Gamma_R$. Lo que sí tiene, y
$\Gamma_R$ no, es una propiedad exacta: vale cero si y solo si el grafo está
desconectado. La etiqueta de régimen describe
entonces una consigna de diseño experimental. No mide la conectividad que cada
mundo acaba teniendo.

% ------------------------------------------------------------------
% Q. Preguntas por nivel (procedia del cuerpo; destino: conclusiones)
% ------------------------------------------------------------------

% ------------------------------------------------------------------
% destino: capítulo de conclusiones
% ------------------------------------------------------------------

